% !TITLE 南北朝
% !SUMMARY 声律的孕育与南北交融
% !ORDER 0

\begin{intro}
南北朝近一百七十年，是中国诗歌从古朴走向精致的关键过渡期。

南朝承袭魏晋遗风，在齐梁年间出现了两个重要的文学运动：一是永明体，沈约、谢朓等人开始有意识地讲究声律，提出"四声八病"之说，为唐代近体诗的格律奠定了基础；二是宫体诗，萧纲、萧绎兄弟把诗歌题材引向闺阁儿女、风花雪月，虽然历来评价不高，却在艺术技巧上做出了精细的探索。谢朓的山水诗清新流丽，"余霞散成绮，澄江静如练"，已经颇有盛唐气象。

北朝则呈现出完全不同的面貌。北方少数民族入主中原，带来了刚健质朴的气质。《敕勒歌》"天苍苍，野茫茫，风吹草低见牛羊"，短短二十七字，却写出了一片苍茫辽阔的天地。北朝民歌中的《木兰诗》，更是叙事诗的杰作——一个女子代父从军的故事，既有英雄气概，又不失女儿情怀。

南北朝后期，庾信由南入北，把南朝的精致技巧与北朝的悲凉气质融为一体，写出了"枯木期填海，青山望断河"这样沉痛的诗句。这种南北交融，正是此后隋唐统一帝国文化整合的预演。没有南北朝在声律、对仗、修辞上的长期积累，就没有唐诗的辉煌。
\end{intro}
